\documentclass[12pt]{article}
\usepackage[utf8]{inputenc}
\usepackage[T1]{fontenc}
\usepackage{hyperref}
\title{Svan Language Resources for K-12 Curriculum Development}
\author{Compiled by Claude Sonnet 5.5 (claude-sonnet-5-5) for Andr\'as Kornai}
\date{}
\begin{document}
\maketitle

\section{Introduction}
Svan (ISO 639-3 sva; Glottocode svan1243) is a Kartvelian language spoken mainly in the Mestia and Lentekhi districts of Svaneti, Georgia. Speaker numbers vary by source: the task description says about 30,000, and Wikipedia lists 14,000 (2015) and classes it as Definitely Endangered. Searches found a grammar and dictionaries but no primer or school material.

\section{Primary Educational Resources}
No primer, textbook or school curriculum was found in the search results.

\section{Linguistic and Reference Materials}
\begin{itemize}
\item Wikipedia entry on Svan: \url{https://en.wikipedia.org/wiki/Svan_language}
\item Tuite, short descriptive grammar of Svan (1998), source record: \url{https://phoible.org/sources/tuite1998}
\item Review of Tuite's Svan grammar by G. Hewitt: \url{https://georgehewitt.net/articles/reviews/141-review-of-kevin-tuite-svan}
\item Article describing Svan grammar and lexis (Kalbotyra, Vilnius University): \url{https://www.zurnalai.vu.lt/kalbotyra/article/download/22211/21382/39845}
\item Georgian--Megrelian--Laz--Svan--English dictionary (Tbilisi State University): \url{https://mail1.tsu.ge/data/image_db_innova/GEORGIAN-MEGRELIAN-LAZ-SVAN ENGLISH DICTIONARY.pdf}
\end{itemize}

\section{Religious and Community Resources}
None found in the search results.

\section{Digital and Audio Resources}
No audio, video or corpus resources were found in these searches.

\section{Curriculum Development Context}
Svan has been written mainly in Georgian script by scholars, and sources conflict on how established the written use is. Curriculum work should involve Svan speakers in Svaneti and Georgian education authorities. Additional searching in Georgian is advisable.
\end{document}